%% kapa: exam=1 round=27 year=2016 period=2 form=A subject=accounting subject_ko=회계학 seq=5 number=45
%% source: Q-Net 감정평가사 기출문제 (cst00309) artlSeq=5169402 file=extracted/1008/5169402/제27회 감정평가사 1차 문제/제27회 감정평가사 2교시문제(A형).pdf
%% conversion: kordoc 4.24.51 -> tools/gen1.py (auto); review=proofread-sonnet; flags=none
\begin{question}{45}
투자부동산의 계정대체와 평가에 관한 설명으로 옳지 않은 것은?
\begin{choices}
  \item 투자부동산을 원가모형으로 평가하는 경우에는 투자부동산, 자가사용부동산, 재고자산 사이에 대체가 발생할 때에 대체 전 자산의 장부금액을 승계한다.
  \item 자가사용부동산을 공정가치로 평가하는 투자부동산으로 대체하는 경우, 사용목적 변경시점까지 그 부동산을 감가상각하고 발생한 손상차손을 인식한다.
  \item 재고자산을 공정가치로 평가하는 투자부동산으로 대체하는 경우, 재고자산의 장부금액과 대체시점의 공정가치의 차액은 기타포괄손익으로 인식한다.
  \item 공정가치로 평가하게 될 자가건설 투자부동산의 건설이나 개발이 완료되면 해당 일의 공정가치와 기존 장부금액의 차액은 당기손익으로 인식한다.
  \item 공정가치로 평가한 투자부동산을 자가사용부동산이나 재고자산으로 대체하는 경우, 후속적인 회계를 위한 간주원가는 사용목적 변경시점의 공정가치가 된다.
\end{choices}
\end{question}
\answer{3}
